\section{Corridor Identification Algorithm}

\label{sec:corridors}
% ══════════════════════════════════════════════════════════════════════════════

Given the predicted habitat suitability map $\mathbf{p}^{(k)} \in
[0,1]^N$ and learned edge weights $\{w_{ij}\}$, we identify migration
corridors through the following procedure.

\subsection{Habitat Patch Detection}

We threshold the suitability map at $\tau = 0.5$ and extract connected
components as habitat patches $\{P_1, \ldots, P_M\}$. Patches smaller
than a species-specific minimum area requirement are removed.

\subsection{Inter-Patch Connectivity}

For each pair of patches $(P_a, P_b)$, we compute the effective
resistance $R_{ab}$ using the learned edge weights as conductances:
\begin{equation}
  R_{ab} = (\mathbf{e}_a - \mathbf{e}_b)^\top \mathbf{L}^{+}
  (\mathbf{e}_a - \mathbf{e}_b)
\end{equation}
where $\mathbf{L}^{+}$ is the pseudoinverse of the graph Laplacian and
$\mathbf{e}_a, \mathbf{e}_b$ are indicator vectors for representative
nodes in each patch.

\subsection{Optimal Corridor Placement}

We formulate corridor design as a minimum-cost flow problem on the
spatial graph. Given a budget of $B$ cells that can be restored to
suitable habitat, we solve:
\begin{equation}
  \min_{\mathbf{z} \in \{0,1\}^N} \sum_{(a,b)} R_{ab}(\mathbf{z})
  \quad \text{s.t.} \quad \sum_i z_i \leq B, \quad z_i = 0 \;
  \forall\, i \in \bigcup_m P_m
\end{equation}
where $z_i = 1$ indicates cell $i$ is selected for restoration. We
approximate this NP-hard problem using a greedy algorithm with
submodular optimization guarantees \citep{nemhauser1978analysis}.

\begin{algorithm}[H]
\caption{Greedy Corridor Placement}
\label{alg:corridor}
\begin{algorithmic}[1]
\REQUIRE Spatial graph $G$, patches $\{P_m\}$, budget $B$
\ENSURE Corridor cells $C \subseteq V$
\STATE $C \leftarrow \emptyset$
\FOR{$b = 1$ to $B$}
  \STATE $v^* \leftarrow \arg\max_{v \notin C \cup \bigcup P_m}
    \Delta R(v \mid C)$
  \STATE $C \leftarrow C \cup \{v^*\}$
\ENDFOR
\RETURN $C$
\end{algorithmic}
\end{algorithm}

Here $\Delta R(v \mid C)$ is the marginal reduction in total effective
resistance from adding cell $v$ to the current corridor set $C$.


% ══════════════════════════════════════════════════════════════════════════════
